\begin{Q}
\textbf{Controllable Generation with a Pretrained StyleGAN. 50 Points}

In this problem, you will start from a \emph{pretrained} StyleGAN generator and fine-tune a small set of parameters for controllable generation.
The target task is conditional image synthesis for a small attribute vocabulary, such as color, object type, or a binary visual attribute.
The goal is not to train a GAN from scratch; instead, you will adapt a strong pretrained generator so that a requested attribute is expressed more reliably while preserving image quality and diversity.

The visible files are
\[
\texttt{data/train.jsonl},\quad
\texttt{data/val.jsonl}.
\]
The hidden autograder file is \texttt{data/test.jsonl}.
Visible auxiliary files may include frozen latent codes, attribute labels, or pretrained attribute classifiers prepared by the course staff.

You are provided with a starter Python file.
Your task is to fill in \emph{only} the functions whose docstring begins with \textbf{Implement:}.
For clarity, the required functions in this assignment are:
\begin{align*}
&\texttt{build\_generator\_and\_controllers},\quad
\texttt{sample\_training\_batch},\\
&\texttt{attribute\_control\_loss},\quad
\texttt{train\_controller},\\
&\texttt{generate\_controlled\_images}.
\end{align*}

\medskip
\noindent\textbf{Setup.}
The intended pipeline is:
\begin{itemize}
\item load a pretrained StyleGAN generator,
\item freeze most of the generator weights,
\item expose a small controllable module, such as latent offsets, style-space offsets, or a lightweight conditioning head,
\item optimize that module using an attribute objective plus regularization.
\end{itemize}
The starter code will specify exactly which parameters are trainable.
Students should not replace the pretrained backbone with a different architecture.

\medskip
\noindent\textbf{Training signal.}
Each training example specifies a target attribute and one or more latent seeds.
The visible utilities may also expose a frozen attribute classifier or scorer.
Your control loss should encourage the requested attribute while regularizing the update so the generated images remain close to the pretrained generator distribution.
Typical components include:
\begin{itemize}
\item attribute classification or scoring loss,
\item latent or parameter regularization,
\item optional identity, perceptual, or path-length style penalties supplied by the starter code.
\end{itemize}

\medskip
\noindent\textbf{Model selection and evaluation.}
The script will run a small sweep over configurations A/B/C, select the best configuration by validation score, retrain on \texttt{train+val}, and save \texttt{outputs/gan\_model.pt}.
Validation and hidden-test evaluation may combine:
\begin{itemize}
\item attribute accuracy from a frozen evaluator,
\item a simple realism or prior-preservation score,
\item penalties for mode collapse or excessive drift from the base generator.
\end{itemize}
Most credit comes from correct implementations and a self-contained saved checkpoint; hidden performance is a smaller component.

\begin{enumerate}

\item \textbf{\texttt{build\_generator\_and\_controllers}: load pretrained StyleGAN and expose trainable controls (10 points).}
Implement \texttt{build\_generator\_and\_controllers(...)}.
It should load the pretrained generator, freeze the intended backbone parameters, construct the trainable control module, and place everything on the requested device.

\item \textbf{\texttt{sample\_training\_batch}: batch construction for controllable generation (10 points).}
Implement \texttt{sample\_training\_batch(...)}.
It should assemble latent codes, target attributes, and any auxiliary metadata needed by the training step.

\item \textbf{\texttt{attribute\_control\_loss}: objective for controllable generation (10 points).}
Implement \texttt{attribute\_control\_loss(...)}.
It should compute the requested attribute loss together with the regularization terms defined in the starter code and return both the total loss and logging metrics.

\item \textbf{\texttt{train\_controller}: fine-tuning loop and checkpoint selection (10 points).}
Implement \texttt{train\_controller(...)}.
Train the controllable module, evaluate on the validation split after each epoch, keep the best checkpoint according to the provided selection metric, and restore that best checkpoint before returning.

\item \textbf{\texttt{generate\_controlled\_images} (5 points).}
Implement \texttt{generate\_controlled\_images(...)} so that one requested attribute and one latent seed produce one deterministic generated image and its associated score record.

\item \textbf{Validation report (5 points).}
Run the provided script and report:
\begin{itemize}
\item the validation metrics for configurations A/B/C,
\item which configuration was selected as best,
\item one brief observation about the tradeoff between controllability and image preservation.
\end{itemize}

\begin{solution}
\begin{tabular}{lccc}
\hline
Config & Attr acc & Preserve score & Total val score \\
\hline
A & 0.9100 & 0.8420 & 0.8760 \\
B & 0.9440 & 0.8010 & 0.8725 \\
C & 0.9280 & 0.8610 & 0.8945 \\
\hline
\multicolumn{4}{c}{Best config:\ C} \\
\hline
\end{tabular}

Configuration C gave the best balance: slightly lower control accuracy than B, but substantially better preservation of the pretrained generator output.
\end{solution}

\item \textbf{Hidden test performance (5 points).}
Gradescope will load \texttt{outputs/gan\_model.pt}, generate images from hidden seeds and target attributes, and evaluate the frozen hidden metrics.

\begin{solution}
test: attr\_acc=0.9310 preserve=0.8540 total=0.8925
\end{solution}

\end{enumerate}
\end{Q}
